\documentclass{situationist-preprint}

\title{Merit in Practice}
\msid{POL-19-b}
\msfamily{POL-19}
\msdiscipline{Political Philosophy}
\mstype{Consequence}
\mskeywords{meritocracy; desert; luck; selection; pay; incentives}

\begin{document}
\maketitle

\begin{abstract}
POL-19-a argues that meritocracy is unjust because merit is partly a product of luck and luck is not deserved, while leaving intact the practice of selecting people for positions by their qualifications. We derive what follows for three practices: selection for positions, the rewards attached to positions, and the public justification of both. Selection by qualifications remains justified, on grounds of efficiency and fair opportunity rather than desert. Differences in reward attached to positions need a justification that does not appeal to desert, such as the incentive argument of the difference principle, and are limited by it. And institutions should not describe the distribution they produce as deserved.
\end{abstract}

\section{Three consequences}

The central result of POL-19-a is Theorem~3.2: in every actual society a meritocratic distribution is not a distribution according to desert, because by Lemma~3.1 it gives different rewards to persons who agree in everything under their control. Corollary~3.3 adds that merit-based selection is not thereby shown to be unjust. Practices in which merit figures can be sorted by which of these results bears on them.

\begin{proposition}\label{p:selection}
Selecting people for positions by their qualifications is justified, where it is, by the efficient performance of the positions and by fair equality of opportunity, not by desert.
\end{proposition}

\begin{proof}
A position exists to be performed, and those who will perform it best are identified, imperfectly, by their qualifications. That justification refers to the position's purpose, not to what the applicant deserves, and it survives Theorem~3.2 of POL-19-a. The principles of POL-01 add the requirement that positions be open to all under conditions of fair equality of opportunity, which limits the ways in which qualifications may be acquired and tested \citep[\S14]{Rawls1971}.
\end{proof}

Two practical consequences follow. A selection procedure may be revised whenever another would better identify those who will perform well, without any loss of what applicants deserve, since they deserved no particular procedure. And preferences that track no qualification, such as preferences for the children of alumni, lack even the justification that qualifications have.

\begin{proposition}\label{p:rewards}
Differences in the rewards attached to positions cannot be justified by the desert of those who hold them; where they are justified, the justification is that they improve the position of the worst-off.
\end{proposition}

\begin{proof}
By Principle~P2 of POL-19-a, persons who agree in what they control deserve the same, so desert cannot justify the difference in pay between two equally conscientious persons, one of whom happens to have capacities the market rewards \citep{Frank2016}. Effort is no exception, since the capacity for effort is itself shaped by endowment and upbringing \citep[\S12]{Rawls1971}. Under the principles of POL-01, a difference in reward is just when it is part of an arrangement that makes the worst-off position better than it would otherwise be, as incentives can do, and to that extent only.
\end{proof}

\begin{proposition}\label{p:rhetoric}
Institutions should not justify the distribution of advantages they produce by the claim that it is deserved.
\end{proposition}

\begin{proof}
By Theorem~3.2 of POL-19-a the claim is false. An institution that makes it misrepresents its own justification, and does so in a way that has costs: those who fare well are encouraged to regard their advantages as earned and those who fare badly to regard their position as their own fault, attitudes that \citet{Young1958} predicted and \citet{Sandel2020} documents.
\end{proof}

\begin{objection}
Without desert, the successful have no claim to what they earn, and the result licenses unlimited redistribution.
\end{objection}

\begin{reply}
They have the claims that just institutions give them: legitimate expectations to what those institutions assign, which POL-19-a does not deny and Proposition~\ref{p:rewards} respects. What they lack is a pre-institutional claim of desert that would limit what just institutions may assign. Redistribution is limited by the principles of justice themselves, including the difference principle's attention to incentives.
\end{reply}

\section{Conclusion}

In practice, the injustice of meritocracy changes the justification of merit-based institutions more than their shape. Selection by qualifications stands; rewards are justified by their effects on the worst-off; and the language of desert is withdrawn from public justification. This is what the conclusion of POL-19, that merit is partly a product of luck and luck is not deserved, requires of the practices that use merit.

\begin{thebibliography}{9}\small
\bibitem[Frank(2016)]{Frank2016} Frank, R. H. (2016). \emph{Success and Luck: Good Fortune and the Myth of Meritocracy}. Princeton University Press.
\bibitem[Rawls(1971)]{Rawls1971} Rawls, J. (1971). \emph{A Theory of Justice}. Harvard University Press.
\bibitem[Sandel(2020)]{Sandel2020} Sandel, M. J. (2020). \emph{The Tyranny of Merit: What's Become of the Common Good?} Farrar, Straus and Giroux.
\bibitem[Young(1958)]{Young1958} Young, M. (1958). \emph{The Rise of the Meritocracy, 1870--2033}. Thames and Hudson.
\end{thebibliography}

\end{document}
